\section{All-contact motion on the original regular source}
\label{sec:v78-motion}
The endpoint field controls more than the two endpoints. Its effect
on every contact decays with distance from the nearer endpoint,
without imposing a positive lower bound on the internal incidences.
This observation permits a uniform estimate for the original smooth
guard, rather than a replacement by protected words.

Keep $q_0=47/53$, $d_j=\min(j,m-j)$ and
$d'_j=\min(j,m-1-j)$. All endpoint derivatives below use
$q=(\alpha_0,\alpha_m)$; internal contact displacements use
arclength. Constants depend only on the fixed physical family and,
where indicated, the originally chosen function $\Theta$.

\begin{lemma}[An incidence-normalized contact inverse]
\label{lem:v78-inverse}
For the optical matrix $A$ of the retained recursion, put
$D_c=\operatorname{diag}(c_1,\ldots,c_{m-1})$ and
$C_A=R_+/[2(1-q_0)]$. For $1\le i,j<m$,
\begin{equation}\label{eq:v78-inverse}
 0\le(A^{-1})_{ij}\le C_A\min(c_i,c_j)q_0^{|i-j|},
 \qquad |(D_c^{-1}A^{-1})_{ij}|\le C_Aq_0^{|i-j|}.
\end{equation}
Consequently, at every regular physical word, its contact positions
$Q_j$ satisfy
\begin{equation}\label{eq:v78-contact-derivative}
 \|D_qQ_j\|\le C_Qq_0^{d_j},\qquad 0\le j\le m,
 \qquad C_Q=R_++\frac{2C_AR_+}{\tau_-q_0}.
\end{equation}
For $m=1$ the internal assertion is empty and the endpoint estimate
still holds. These are derivatives at actual regular trajectories;
no common neighborhood of physical continuation is asserted.
\end{lemma}
\begin{proof}
Write $A=D(I-P)$, with $D$ its diagonal. The entries of $P$ are
nonnegative, vanish except at neighboring indices, and have row sums
at most
\[
 \frac{2/\tau_-}{2/\tau_-+2/R_+}=q_0.
\]
Also $D_{jj}^{-1}\le Rc_j/2$. In the Neumann series for
$(I-P)^{-1}D^{-1}$, the $(i,j)$ entry vanishes until the power
$|i-j|$. Summing the remaining geometric series gives the bound
with $c_j$. Symmetry gives the bound with $c_i$, and division by
$c_i>0$ proves the second inequality. This is the same normalized
inverse underlying the retained tapered-contact construction; here
only its first derivative is needed, at arbitrary regular endpoints.

The internal action Hessian in oriented arclengths is $D_cAD_c$.
After division of its stationary equations by the internal cosines,
endpoint forcing has only the first and last entries, of absolute
sizes $t_0c_0\,du_0$ and $t_{m-1}c_m\,du_m$. Tangent signs may
change these signs but not their sizes. Thus the displacement of
internal arclength $s_j$ obeys
\[
 |ds_j|\le\frac{C_A}{\tau_-}
     \bigl(q_0^{j-1}|du_0|+q_0^{m-1-j}|du_m|\bigr).
\]
Both exponents are at least $d_j-1$, and
$|du_i|=R|d\alpha_i|$. The displayed choice of $C_Q$ bounds the
resulting Euclidean operator norm and the two endpoint derivatives.
The Euclidean speed of a circle parameterized by arclength is one,
which proves \eqref{eq:v78-contact-derivative}.
\end{proof}

\begin{theorem}[Uniform motion of every physical margin]
\label{thm:v78-margins}
There is $C_D<\infty$, independent of $m,R,n,k$ and of every
positive incidence and clearance, with the following property.
Along any regular endpoint curve of length $\ell_q$ in the exact
event and in one physical word,
\begin{align}
 |c_j(x)-c_j(y)|+|p_j(x)-p_j(y)|
      +\operatorname{dist}_{\mathbb S^1}(\alpha_j(x),\alpha_j(y))
       &\le C_Dq_0^{d_j}\ell_q,\label{eq:v78-contact-motion}\\
 |\Delta_j(x)-\Delta_j(y)|&\le C_Dq_0^{d'_j}\ell_q.
                                      \label{eq:v78-clearance-motion}
\end{align}
In particular, for a forward survivor of the field in
Theorem~\ref{thm:v77-field}, the right sides are bounded by,
respectively, $C_D C_Xq_0^{d_j}\sqrt h$ and
$C_D C_Xq_0^{d'_j}\sqrt h$, where
$C_X=\sqrt{2\Lambda_\circ}/\lambda_\circ$.
\end{theorem}
\begin{proof}
Differentiate a flight direction:
\[
 dv_j=\ell_j^{-1}(I-v_jv_j^{\mathsf T})(dQ_{j+1}-dQ_j).
\]
The depths of its two endpoints are at least $d'_j$. Its derivative
is consequently bounded by $2C_Qq_0^{d'_j}/\tau_-$.
The normal and tangent derivatives at a contact are bounded by
$C_Qq_0^{d_j}/R_-$. Express $c_j$ and $p_j$ as scalar products
of the appropriate unit flight velocity with these vectors. At the
terminal contact use the incoming velocity, with the sign giving
positive incidence; its tangential component is unchanged by
reflection. Since adjacent depths differ by at most one, these
formulas prove the first estimate with a fixed $C_D$. In particular
we do not differentiate $\sqrt{1-p_j^2}$ by dividing by $c_j$.

The Hausdorff displacement of a segment is no greater than the
larger displacement of its two endpoints. Distance from that segment
to a fixed third disk is a one-Lipschitz function of this Hausdorff
displacement. Taking the minimum over third disks preserves the
same bound. The minimizing disk may change; differentiability of
that minimum is not required. The excluded incident disks are fixed
on the word. Integration along the curve proves
\eqref{eq:v78-clearance-motion}. Finally use the endpoint-curve
length estimate of Theorem~\ref{thm:v77-field}.
\end{proof}

\begin{corollary}[The unchanged guard is uniformly controlled]
\label{cor:v78-guard}
For every forward survivor in the original exact event,
\begin{equation}\label{eq:v78-guard}
 |W^\varepsilon(\Phi_hx)-W^\varepsilon(x)|
   \le\min\{1,K_\Theta\varepsilon^{-1}\sqrt h\},\qquad
 K_\Theta=\frac{4\operatorname{Lip}(\Theta)C_DC_X}
                         {1-q_0^{3/4}}.
\end{equation}
No incidence cutoff, clearance cutoff, section-distance factor or
monotonicity assumption on $\Theta$ is added.
\end{corollary}
\begin{proof}
For factors in $[0,1]$, the difference of their products is bounded
by the sum of their absolute differences. A contact factor in
\eqref{eq:v76-guard} changes by at most
$\operatorname{Lip}(\Theta)C_DC_X\varepsilon^{-1}
q_0^{3d_j/4}\sqrt h$; the analogous flight factor has depth $d'_j$.
There are at most two contacts and two flights at any depth. Their
geometric series gives \eqref{eq:v78-guard}. Both weights belong
to $[0,1]$, proving the additional bound by one. The argument uses
only the Lipschitz constant and range of the original smooth function.
\end{proof}

\begin{proposition}[Where a finite forward exit can occur]
\label{prop:v78-exit-localization}
Let $s_j$ be distance from the collision state $(\alpha_j,p_j)$
to the boundary of $Y_R^*$ in the product metric. There is a common
$C_E$ such that any source point whose forward lifetime in the
regular exact event is at most $h$ belongs to the union
\begin{align}
 \mathcal B_h={}&\bigcup_{j=0}^m
     \{c_j\le C_Eq_0^{d_j}\sqrt h\}
 \ \cup\ \bigcup_{j=0}^{m-1}
     \{\Delta_j\le C_Eq_0^{d'_j}\sqrt h\}\notag\\
 &\hspace{10mm}\cup\ \bigcup_{j=0}^m
     \{s_j\le C_Eq_0^{d_j}\sqrt h\}.\label{eq:v78-exit-layer}
\end{align}
Consequently its exit mass is bounded above by
$\int W^\varepsilon\one_{\mathcal B_h}\,d\mu_t$.
This is a source inclusion, not a pointwise concentration estimate.
\end{proposition}
\begin{proof}
The finite endpoint-curve length gives limiting contact data up to
a finite lifetime. A forward trajectory cannot reach an interior
critical point, by Theorem~\ref{thm:v77-field}. If all limiting
incidences, clearances and distances to section edges were positive,
the regular reflection equations, strict first-hit tests and locally
constant section decisions would continue the same physical word
and exact event. The endpoint Jacobian remains nonzero there, so
the roof field also continues, a contradiction. Thus some listed
margin tends to zero. Apply Theorem~\ref{thm:v78-margins}, including
the one-Lipschitz property of distance to a closed edge set. Taking
the union over the possible margins and retaining the original
weight proves the mass bound. A section edge is charged whether
or not another simultaneous change might preserve the occupation;
this can enlarge the upper bound but cannot omit a true exit.
\end{proof}
